% =====================================================================
% Section: The space of computations -- spacetime from non-computability
% Drop-in for Paper 1 (or Paper 2; this is the GR/spacetime arc). Requires:
%   \usepackage{amsmath, amssymb}
%   \usepackage[backend=biber, style=phys]{biblatex}   % or your chosen style
% Internal cross-references (\ref{...}) are placeholders; repoint to your labels:
%     sec:complex-amplitudes -- self-reference -> negative probability ->
%                               complex (imaginary) information
%     sec:postulates         -- the quantum-postulate derivation
%     sec:fanout-boundary    -- the FANOUT boundary
%     sec:gauge-group        -- the division-algebra -> G_SM derivation
% Citation keys refer to the companion .bib; reconcile with your Mendeley keys.
% British English throughout.
% =====================================================================

\section{Spacetime from non-computability}
\label{sec:space-of-computations}

We'll now make a brief diversion to take up this idea that non-computable steps can be measured by imaginary information \ref{sec:measuring-non-computability}.  Since Spencer-Brown assures us that Bertrand Russell himself though well of the idea,\autocite{Spence-Brown-1969} it's not completely crazy.

The programmes that derive spacetime from quantum mechanics---Van~Raamsdonk's
\textquote{building spacetime with entanglement} \autocite{VanRaamsdonk2010}, the
ER\,$=$\,EPR conjecture \autocite{MaldacenaSusskind2013}, and M\"uller's
information-theoretic recovery of three-dimensional space from the qubit
\autocite{MuellerMasanes2013}---take entanglement or information as primitive and
grow geometry from it. The route taken here is the same in spirit but inverted in
order: rather than positing spatial dimensions and forbidding superluminal
signalling, we \emph{define} spatial separation as a measure of non-computability
and read the dimension, the signature, and the Lorentz group off the result. The
proposal originates in the author's CASYS'05 paper \autocite{small_2006}, which
first suggested treating spacetime as a space of computations and attempted to
justify three dimensions on this basis; the present account tightens and extends
it. This is the framework's most structural and least formally complete arc; the
three-dimensionality result is a structural conclusion awaiting a formal proof,
and several of the links below are conjectural readings. Each is flagged as such,
so the reader sees exactly where the chain is load-bearing and where it is an
invitation.

\subsection{Timelike and spacelike distance as computation}
\label{sec:comp-distance}

A computation carries a state from axioms to a theorem in a number of steps.
Counting the computable steps between two states gives a real measure---time. A
physical clock counts completed irreversible loops, so proper time literally
counts loop completions and the metric $\mathrm{d}\tau^{2}=g_{\mu\nu}\,
\mathrm{d}x^{\mu}\mathrm{d}x^{\nu}$ is a loop-density tensor, the conversion from
coordinate intervals to loop counts; Toffoli's reading of the Lagrangian as a
count of action quanta is the same statement \autocite{toffoli_2003}.

Spacelike separation is the complementary, non-computable measure. Define the
non-computable distance between two propositions, relative to an axiom system, as
the number of bits that must be added to the axioms to make the second derivable.
A closed timelike curve is a computation whose output is its own input; by the
fixed-point theorem it has a unique consistent solution or none, and the
existence of that solution is one bit not derivable from any axiom system that
does not already contain the loop---one bit per curve, in the manner of Chaitin's
$\Omega$ \autocite{Chaitin1998}. The spacelike distance between two events is then
the number of such curves that would have to be closed to make them causally
connected. Because each added bit is a non-computable proposition, it is carried
by the imaginary part of the complex information measure
$S=-\ln\lvert p\rvert-\mathrm{i}\theta$ derived in
Section~\ref{sec:complex-amplitudes}: computable distance is counted in real
units, non-computable distance in imaginary ones. The Minkowski signature, with
spacelike intervals imaginary relative to timelike ones, is the immediate reading.
That the resulting structure is a genuine metric space, and that the signature
follows rather than being fitted, are open problems---at present the
identification is motivated by reproducing what nature does, not proved.

\subsection{Three dimensions of space}
\label{sec:three-dimensions}

The closed timelike curves that define spacelike distance require information to
be parallel transported consistently around a loop, and consistent parallel
transport on a sphere requires it to be parallelisable. By Adams' theorem the
only parallelisable spheres are $S^{1}$, $S^{3}$ and $S^{7}$
\autocite{adams_1960}, the unit spheres of $\mathbb{C}$, $\mathbb{H}$ and
$\mathbb{O}$. Impose now the requirement that information arriving from outside
the causal zone must appear to possess its own causal history---that is, to be the
result of a sequence of information-preserving operations, hence a sequence of
group operations. Rotations of $S^{1}$ and $S^{3}$ form the groups $U(1)$ and
$SU(2)$ and preserve information; rotations of $S^{7}$ are non-associative and do
not form a group, so $S^{7}$ cannot serve as a ruler. The same non-associativity
forbids a metric outright: path composition $d(A,C)$ through $B$ depends on the
bracketing, so the triangle inequality is not merely violated but ill-defined, and
no geodesic structure exists. Since the abelian $S^{1}$ alone cannot carry
non-abelian structure, spacelike distance must be three-dimensional. Dimension
three is an output, not an input \autocite{small_2006}.

The division of labour is then exact. $S^{1}$ ($\mathbb{C}$) contributes the time
dimension and the $U(1)$ of electromagnetism; $S^{3}$ ($\mathbb{H}$) contributes
the three spatial dimensions and the $SU(2)$ of the weak interaction; $S^{7}$
($\mathbb{O}$) contributes no geometry but, as a filter on what information may be
exchanged, the $SU(3)$ of the strong interaction. A sphere that cannot be a ruler
can still be a filter, which is why the strong force resists geometric
interpretation while electromagnetism and the weak force do not. This is the same
$\mathbb{C}\oplus\mathbb{C}^{3}$ split that yields $G_{\mathrm{SM}}$ in
Section~\ref{sec:gauge-group}, now read as the reason spacetime is $(3+1)$-%
dimensional. The conclusion converges with M\"uller and Masanes, who prove from
information-theoretic postulates, assuming no quantum theory, that minimal
direction-carrying systems with continuous reversible dynamics force $d=3$ and
quantum theory on two qubits \autocite{MuellerMasanes2013}, and with the result
that entanglement exists only for a three-dimensional Bloch ball
\autocite{MasanesMuellerEtAl2014}; the framework supplies the reason such systems
exist, Cayley--Dickson doubling from self-reference, while their work supplies
the rigorous dimensional uniqueness. The status of the framework's own route is a
structural result requiring formal proof: the one-bit-per-curve identification and
the precise reduction of the $S^{7}$ filter to $SU(3)$ rather than its full
automorphism group $G_{2}$ are named open problems.

\subsection{Recovering the Lorentz group}
\label{sec:lorentz}

A measurement changes what the observer knows, hence, relationally, the causal
structure: it tips the light cone. This is Penrose's association of cone-tipping
with state reduction \autocite{Penrose1996} run in the opposite causal direction,
reduction driving the tipping rather than the reverse, and the two are then one
operation in two charts---both described by $SL(2,\mathbb{C})$. The algebraic
backbone is the division-algebra correspondence
$SL(2,\mathbb{A})\cong\mathrm{Spin}(1,d{+}1)$ \autocite{baez_2002}: $\mathbb{R}$
gives $1{+}1$, $\mathbb{C}$ gives $3{+}1$, $\mathbb{H}$ gives $5{+}1$ and
$\mathbb{O}$ gives $9{+}1$ dimensions. The computability split that singles out a
preferred $\mathbb{C}$ fixes the fundamental state space as $\mathbb{C}^{2}$ and
hence $SL(2,\mathbb{C})=\mathrm{Spin}(1,3)$: the Lorentz group of $3{+}1$
spacetime. The qubit and Minkowski space are not merely both Euclidean-and-three%
-dimensional by coincidence; the Hermitian $2\times2$ matrices are
$\mathbb{R}^{1,3}$, $SL(2,\mathbb{C})$ acts on them as the Lorentz group, and the
qubit density operators fill the future cone---the correspondence M\"uller's group
develops operationally, recovering Lorentz transformations from quantum
communication \autocite{HohnMueller2016} and the complex and quaternionic qubit
from the relativity of simultaneity \autocite{GarnerMuellerDahlsten2017}. The
$SL(2,\mathbb{A})$ correspondence is rigorous; its identification with
measurement-as-tipping is the framework's structural reading.

\subsection{The relationalist commitment}
\label{sec:relationalism}

General relativity is relational: by the hole argument, spacetime points carry no
intrinsic identity, and the metric is fixed by the network of coincidences and
interactions, not by a pre-existing stage \autocite{Rovelli1996a}. The
computational picture is relational by construction and explains why it must be:
a background metric would require spatial distances to be pre-computable, and
non-computability forbids exactly that. General covariance is the statement that
distances have no pre-assigned values and depend on what information is actually
present. The metric is therefore dynamical not by stipulation but because it
measures a quantity that cannot be fixed in advance.

\subsection{Jacobson, holography, and imaginary information}
\label{sec:jacobson}

Jacobson derives the Einstein equation as an equation of state, from the Clausius
relation $\delta Q=T\,\mathrm{d}S$ across every local Rindler horizon together
with the Bekenstein proportionality of entropy to horizon area
\autocite{Jacobson1995, Bekenstein1973}. In the computational reading this is the
self-consistency condition on the loop-counting network: the Einstein equation is
the requirement that the bookkeeping of computation be internally consistent
across surfaces. The holographic principle, that the information in a region is
bounded by its bounding area \autocite{tHooft1993, Susskind1995}, then acquires a
natural reading here. A surface area is a squared distance; if spacelike distance
is measured in imaginary units of information, an area is measured in the square of
those units, that is in \emph{negative} information---entropy. The holographic
link between area and entropy thus requires the square root of negative
information, and that square root is precisely the imaginary unit of the complex
amplitude, the observer's self-referential phase. Read this way, Jacobson's
account of information crossing surfaces implicitly assigns an imaginary
information value to spacelike distance, and the $\mathrm{i}$ of the Minkowski
signature is the same $\mathrm{i}$ as the $\mathrm{i}$ of the quantum amplitude.
This identification is the most speculative step in the section: it is
structurally suggestive and it makes the holographic form unsurprising, but it is
a conjectural reading rather than a derivation, and is offered as such.

\subsection{Non-locality already present in general relativity}
\label{sec:gr-nonlocality}

The relationalist commitment has a consequence, set out in the author's 2025
V\"axj\"o poster \autocite{Small2025} and developed here, that the wider
literature has not drawn. If spacetime is the network of interactions, then a system that
has not yet interacted with any participant in that network is not a node in it:
it is, strictly, outside spacetime until it interacts. An entangled pair that has
not been measured is one such system---a single relational object, not two
spatially separated ones. The electron in a sealed bottle illustrates it: the
electron's position lies in the observer's past light cone, but its spin does not
enter the observer's causal past until it is measured, even though the electron is
spatially close. Entanglement is therefore a relation that does not yet live in
spacetime; the spatial separation later assigned to the partners is a derived
fact, established only when the relation is \textsc{fanned out} into the network
(Section~\ref{sec:fanout-boundary}). The upshot is that general relativity, taken
relationally and seriously, already contains quantum non-locality: entanglement is
the pre-spatial relational structure, and what looks like spooky action between
distant points is the resolution of a relation that was never spatially separated
to begin with. This is the same direction of travel as building spacetime from
entanglement \autocite{VanRaamsdonk2010} and as ER\,$=$\,EPR
\autocite{MaldacenaSusskind2013}, but reached from the relational and
computational side rather than from holographic duality, so it stands as a
parallel argument rather than a corollary of those results. The claim that the
non-locality is \emph{already there}, unrecognised, in the relational reading of
general relativity is the framework's interpretive thesis, advanced in
\autocite{small_2006, Small2025} and flagged as a thesis rather than a theorem.

\medskip
\noindent The spacetime arc therefore yields, from one definition of spacelike
distance as non-computability: the signature, three spatial dimensions, the
Lorentz group, the dynamical relational metric, a reading of the holographic
principle, and the presence of non-locality within general relativity. Of these,
the $(3+1)$ count and the Lorentz recovery are structural results that converge
with independent rigorous work and await their own formal proof; the metric-space
status, the signature, the one-bit-per-curve measure, the imaginary-information
reading of holography, and the general-relativistic non-locality are conjectural,
and are named here precisely so that the boundary between what is shown and what
is proposed is visible.

% =====================================================================
% Subsection: Open problems and prospects (closes the spacetime section)
% Append at the end of space_of_computations_2026-06-07.tex, inside
% \section{The space of computations...} (sec:space-of-computations).
% Requires the same preamble. New citation keys are in the updated
% space_of_computations_refs .bib. Internal \ref placeholders:
%     sec:gr-nonlocality  -- the GR non-locality subsection
%     sec:three-dimensions, sec:jacobson  -- as labelled in the main section
%     sec:gauge-group, sec:k3-sterile  -- the SM gauge-group / K3 sections
% British English throughout.
% =====================================================================

\subsection{Open problems and prospects}
\label{sec:spacetime-open-problems}

The spacetime arc is the framework's frontier, and honesty about its edges
matters more here than anywhere else in the paper. We set out the genuine
problems and the prospects together, deferring the deepest to a second paper
rather than papering over them.

\paragraph{The substrate and the generation of events (deferred).}
A computation runs on something; a network of events presupposes that events
occur. Both are the same objection in two forms---the network's hardware and its
clock---and the original CASYS'05 statement of the programme raised the
implementation regress and set it aside \autocite{small_2006}. The framework's
response is that a self-referential structure is the only kind that need not
presuppose a meta-substrate: it generates its own. The natural algebraic emblem
is $E_{8}$, the unique simple Lie group whose adjoint representation is also its
smallest faithful one, so that what it is and what it acts on coincide; the $240$
unit integral octonions that generate it are exactly the $E_{8}$ root system, so
the self-action is welded to the division-algebra spine of
Section~\ref{sec:gauge-group}. This is Wheeler's self-excited circuit
\autocite{Wheeler1990} given an algebraic body, and it converges with Furey's use
of an algebra acting on itself to generate particle content
\autocite{furey2018}. The decisive caveat, which we state rather than hide, is
that the bootstrap is self-consistent but not self-specifying: it does not
eliminate the non-computability identified throughout, it relocates it to a
determined place. The substrate question is, on the framework's own terms,
provably without a computable answer from inside---it is the irreducible residue,
the bump in the carpet that the whole construction has been moving rather than
removing. The leap from ``$E_{8}$ is self-dual'' to ``$E_{8}$ is the substrate''
is at present an emblem awaiting operational content (an $E_{8}$-equivariant
dynamics, or a state space built as a homogeneous space of $E_{8}$), and that
content, together with how a tick is generated at all, is deferred to the second
paper.

\paragraph{Dark matter (deferred).}
If gravity is the bookkeeping of the interaction network and dark matter
interacts only gravitationally, how does it participate in a network it does not
exchange information with? The prospect, to be developed in the second paper, is
that this is a feature rather than a defect. In the division of labour of
Section~\ref{sec:three-dimensions}, participation in the metric---counting loops
in the $S^{1}$ and $S^{3}$ directions---is universal, carried by anything with a
worldline and a tick-rate, whereas participation in the gauge filters, the
$S^{3}$ of the weak interaction and the $S^{7}$ of colour, is selective, engaged
only by charged matter. A node that counts loops but engages no filter
\emph{is} matter that gravitates without otherwise interacting, and the
framework's Class-4 a-chiral sterile neutrino (Section~\ref{sec:k3-sterile}) is
the candidate that realises it. The condition for this to be coherent rather than
circular is that metric participation be established as prior to, and independent
of, gauge participation; making that priority precise is the work deferred.

\paragraph{CP violation, the arrow of time, and baryogenesis.}
CP violation appears to demand a time asymmetry owing nothing to the entropic
arrow, which would seem to require a preferred direction built into spacetime
itself. The framework dissolves the appearance rather than accommodating it. CP
violation needs an irreducible complex phase, the CKM phase $\delta$; by Spekkens'
result that negativity and contextuality are the same notion of nonclassicality
\autocite{Spekkens2004a}, a complex phase is a signature of contextuality; so CP
violation should be contextual, sensitive to the near environment of the decaying
system---which is exactly what Andrianov, Taron and Tarrach established for
neutral kaons in matter, where the environment breaks C, irreversibility breaks
T, and the effective parameters follow a Lindblad equation
\autocite{Andrianov2001}. The framework's reading is that the resulting
informational asymmetry \emph{constitutes} the arrow rather than presupposing it:
CP violation, entropy increase, and the direction of time are one informational
asymmetry seen three ways, not an independent second arrow, so the worry about a
direction built into spacetime does not arise. On this footing the contextual
dependence sets up a positive feedback loop for baryogenesis---an initial
fluctuation makes the environment slightly matter-dominated, contextual
enhancement raises the effective CP violation, more matter is produced, and the
loop runs---with $\delta$ a cosmological fixed point rather than a fundamental
constant. The honest limit of this is that a feedback loop explains amplification
and lock-in, not the sign selection or the baryon-number violation that
Sakharov's conditions still require \autocite{Sakharov1967}; the stronger claim
the framework is positioned to make, that the matter direction and the time
direction are co-selected by the same self-referential asymmetry, needs a source
of baryon-number violation and is left open.

\paragraph{Spatial infinitude and the status of inflation.}
There is no upper bound on non-computability, so the space of computations is
unbounded: the universe is spatially infinite, and being a network without a
boundary, naturally flat. Infinitude on its own, however, settles only flatness,
which was never inflation's principal task. The horizon problem is, and an
infinite space does nothing for it. The framework's actual lever against
inflation is the relational thesis of Section~\ref{sec:gr-nonlocality}: if
pre-spatial entanglement correlates regions before \textsc{fanout} spreads them
into spatial separation, then the observed large-scale homogeneity is fixed
relationally, before any question of causal contact, and the horizon problem
dissolves without an inflaton. This is a real argument, but it carries a real
bill: inflation also generates the nearly scale-invariant spectrum of primordial
fluctuations seen in the microwave background \autocite{Planck2020}, and any
account that dispenses with inflation owes an alternative origin for that
spectrum. That, not the horizon problem, is the hard empirical debt, and it is
named here rather than discharged.

\paragraph{Expansion and acceleration.}
If spacelike distance is the information separating worldlines, then as worldlines
accumulate distinct interaction histories they accumulate information-distance,
and the universe expands. The picture has the right exclusions built in: bound
systems interact continuously, share information, and so keep low
information-distance---they do not expand---while only worldlines that decouple
diverge, recovering expansion on cosmological scales alone. The feedback in which
greater separation means less light-cone overlap and so faster divergence,
$\mathrm{d}d/\mathrm{d}t \propto d$, yields exponential divergence, a de~Sitter
phase, and hence accelerating expansion without a separate dark-energy term
\autocite{Riess1998}. The limit to be honest about is that pure exponential
divergence is the late-time behaviour only; reproducing the full expansion history,
with its radiation- and matter-dominated eras, and fixing the Hubble rate
quantitatively, remains to be done.

\paragraph{The $S^{7}$ sector as internal structure.}
That $S^{7}$ cannot carry a metric does not make it disappear. Its closed
timelike curves carry the same non-computable information structure as those of
$S^{1}$ and $S^{3}$, but lacking associativity they support algebra---gauge
charge---in place of geometry---distance. This is why gauge symmetry has so often
looked like a set of compactified extra dimensions: the Kaluza--Klein and string
intuition is tracking something real in the $S^{7}$ structure, but misdiagnoses
it as spatial. In the framework these degrees of freedom have no size because
they have no metric; they are internal in the exact sense that they enter the
information-exchange that constitutes gauge charge while contributing nothing to
the loop-counting that constitutes distance. The resemblance to rolled-up
dimensions is therefore genuine and instructive, and the diagnosis that they are
algebraic rather than geometric is what separates this account from theories that
add real spatial dimensions to carry the same structure.

\medskip
\noindent These mark the boundary of the programme as it stands. The substrate,
the generation of events, and the network status of dark matter are deferred to
the second paper; the contextual account of CP violation and baryogenesis, the
relational case against inflation, the information-divergence account of
expansion, and the algebraic reading of the $S^{7}$ sector are offered as
prospects, each with its outstanding debt named so that the line between argument
and aspiration stays visible.